%قبل از شروع مسئله ترجیحاً چیزی ننویسید. اما اگر خواستید بنویسید، در قالب لیستی (فقرات و شمارش) بنویسید.
%برای اضافه کردن نمره، بعدی از شروع‌مسئله، [x] را اضافه کنید و به جای x  نمره مورد نظر را قرار دهید. برای زیرمسئله‌ها هم به همین صورت است.
%\شروع‌مسئله[10]
\شروع‌مسئله

% صورت اصلی مسئله را در این‌جا بنویسید.

یک درخت با 
 n  
راس در اختیارتان قرار گرفته است. هر یال از این درخت نیز، یک وزن صحیح مثبت دارد.  الگوریتمی با مرتبه زمانی
  O(n)  
ارائه دهید که طولانی‌ترین مسیر در این درخت را پیدا کند.

\solution{
	

فرض کنید طولانی‌ترین مسیر در درخت، بین دو راس
 A 
و
 B 
باشد. از یک راس تصادفی، الگوریتم
 DFS 
را اجرا می‌کنیم و آن راس را
 x 
می‌نامیم.  با این عمل، راسی که بیشترین فاصله از 
 x 
دارد را پیدا می‌کنیم. این راس، باید یا
 A 
باشد یا
 B 
. این موضوع را می‌توانید با استفاده از برهان خلف ثابت کنید. فرض کنید این راس هیچ‌کدام از 
A
یا
B
نباشد و ببینید که در آن صورت، فرض این که طولانی‌ترین مسیر بین 
 A 
و 
 B 
است دچار اشکال می‌شود.
(منظور از بیشترین فاصله، جمع وزن یال‌های مسیر است. کافی است هنگام پیمایش درخت برای هر راس نگه‌داریم که از کدام راس به آن وارد شده‌ایم و به این ترتیب می‌توان وزن مسیر طی شده از راس مبدا تا هر راس را به دست آورد) 
 
 
 بدون از دست رفتن کلیت مسئله فرض کنید دورترین راس از 
x
راس
A 
باشد. حالا از 
A 
مجدداً 
DFS
را اجرا کنید و دورترین راس را پیدا کنید که 
B
خواهد بود.
دو بار الگوریتم جست‌وجوی عمق اول اجرا شده است که چون گراف یک درخت است، مرتبه زمانی خطی خواهد داشت.

}



%%%%%%%%%%%%%%%%%%%%%%%%%%%%%%%%%%%%%%%%%%%%%%%%%%%%%
% برخی دستورات مفید برای تمرین‌ها. %
%%%%%%%%%%%%%%%%%%%%%%%%%%%%%%%%%%%%%%%%%%%%%%%%%%%%%

% فرمول‌های نمونه:
%$f(x) = \Theta(x)$
%$T(x, y) \leq \log(x) + 2^y + \binom{2n}{n}$
% {\tt code\_function}


% لیست موارد بدون شماره‌گذاری:
%\شروع{فقرات}
%    \فقره مورد اول 
%        \شروع{فقرات}
%            \فقره مورد اول داخلی 
%                \شروع{فقرات}
%                    \فقره اولین مورد خیلی داخلی 
%                    \فقره دومین مورد خیلی داخلی 
%                \پایان{فقرات}
%            \فقره مورد دوم داخلی 
%        \پایان{فقرات}
%    \فقره مورد دوم 
%\پایان{فقرات}

% لیست موارد شماره‌گذاری شده:
%\شروع{شمارش}
%    \فقره یک مورد 
%    \فقره یک مورد دیگر 
%    \فقره باز هم یک مورد دیگر
%\پایان{شمارش}

% برای نوشتن یک فرمول چند خطی که یک شماره بخوره. برای حالت‌های دیگر به سایت زیر مراجعه کنید:
% https://www.overleaf.com/learn/latex/Aligning_equations_with_amsmath
%\begin{align} 
%    \begin{split}
%        (x+y)^3 &= (x+y)^2(x+y) \\
%                &= (x^2+2xy+y^2)(x+y) \\
%                &= (x^3+2x^2y+xy^2) + (x^2y+2xy^2+y^3) \\
%                &= x^3+3x^2y+3xy^2+y^3
%    \end{split}                                 
%\end{align}

% نوشتن ماتریس و آرایه دوبعدی
%\begin{align}
%    A = 
%    \begin{bmatrix}
%        A_{11} & A_{21} \\
%        A_{21} & A_{22}
%    \end{bmatrix}
%\end{align}

% اضافه کردن تصویر با شرح. برای اضافه کردن تصویر بدون شرح، فقط دستور درج‌تصویر را استفاده کنید.
%\شروع{شکل}[h!]
%\تنظیم‌ازوسط
%  \درج‌تصویر[پهنا=0.5\پهنای‌متن]{img.jpg}
%\شرح{شرح تصویر}
%\برچسب{تصویر:۱}
%\پایان{شکل}
%\پایان{زیرمسائل}

%کد زیر را ببینید.
%%\lstinputlisting{source_filename.py}
%\begin{latin}
%\begin{lstlisting}
%# Import the modules
%import sys
%import random
%
%ans = True
%
%while ans:
%question = raw_input("Ask the magic 8 ball a question: (press enter to quit) ")
%
%answers = random.randint(1,8)
%
%if question == "":
%sys.exit()
%
%elif answers == 1:
%print "It is certain"
%
%elif answers == 2:
%print "Outlook good"
%
%elif answers == 3:
%print "You may rely on it"
%
%elif answers == 4:
%print "Ask again later"
%
%elif answers == 5:
%print "Concentrate and ask again"
%
%elif answers == 6:
%print "Reply hazy, try again"
%
%elif answers == 7:
%print "My reply is no"
%
%elif answers == 8:
%print "My sources say no"
%\end{lstlisting}
%\end{latin}
